\chapter{Revisão da Literatura}

Nesta seção são abordados os principais fundamentos e características técnicas das tecnologias Rust e WebAssembly. Inicialmente, apresenta-se a linguagem Rust, sua performance e sua proposta de garantir segurança de memória sem o uso de coletor de lixo, delegando o controle de alocação ao compilador por meio do modelo de \emph{ownership}. Em seguida, discute-se o uso de WebAssembly como padrão binário de baixo nível suportado pelos principais navegadores, abordando sua arquitetura, modelo de execução e o processo de compilação de código Rust para esse formato.

\section{Rust: fundamentos, segurança e performance}

Rust é uma linguagem de programação de sistemas desenvolvida a partir da necessidade de conciliar alto desempenho e segurança no desenvolvimento de software. Desde que a Mozilla começou a patrocinar seu desenvolvimento em 2010, a linguagem foi ganhando espaço na comunidade devido à forma como ela lida com dois aspectos historicamente desafiadores no desenvolvimento de software: o desempenho e a segurança de memória.

\subsection{Segurança: o Modelo de Ownership e Borrowing}

Se o desempenho de Rust já se mostra expressivo, é na segurança que a linguagem realmente se destaca. Estima-se que cerca de 70\% das vulnerabilidades de segurança catalogadas pelo Microsoft Security Response Center têm origem em erros de gerenciamento de memória \cite{bugden2022}. Nesse contexto, o Rust aborda diretamente esse problema por meio de um mecanismo denominado \emph{ownership} (propriedade).

O conceito é baseado em um princípio: cada valor em Rust tem um único dono, responsável por seu gerenciamento. Quando o proprietário deixa de existir ou sai de escopo, o valor associado é automaticamente desalocado da memória. Isso impede, por exemplo, que um programa tente usar um dado que já foi descartado, erro conhecido como \emph{use after free}, comum em C e frequentemente explorado por hackers.

O Algoritmo 1 apresenta um exemplo do funcionamento do \emph{ownership}. Inicialmente, a variável nome é proprietária de uma String. Quando essa variável é atribuída a linguagem, a propriedade do valor é transferida, caracterizando uma operação de \emph{move}. A partir desse momento, nome deixa de ser uma variável válida para acessar o valor, enquanto linguagem passa a ser sua proprietária.

\begin{lstlisting}[
    language=Rust,
    caption={Exemplo de \emph{ownership}: cada valor tem um único dono},
    label={alg:ownership}
]
fn main() {
    let nome = String::from("Rust");
    // O ownership de 'nome' é transferido para 'linguagem'.
    let linguagem = nome;
    // A linha abaixo causaria erro de compilação,
    // porque 'nome' não é mais o dono do valor.
    println!("{}", nome); // erro!
    println!("{}", linguagem); // funciona normalmente
}
\end{lstlisting}
\fontealgoritmo{Os autores}

A tentativa de utilizar nome após a transferência de propriedade resulta em erro de compilação, pois o valor já pertence à variável linguagem. Dessa forma, o compilador impede que uma variável continue sendo utilizada como se ainda fosse proprietária de um recurso que já foi transferido.

Quando há necessidade de compartilhar um valor entre diferentes partes do programa sem transferir o \emph{ownership}, existe o \emph{borrowing} (empréstimo). Nesse mecanismo, uma parte do código pode acessar temporariamente um valor pertencente a outra parte sem assumir sua propriedade. As regras são: ou várias partes do código leem o dado ao mesmo tempo, referências imutáveis, ou uma única parte pode modificá-lo, referência mutável. As duas situações nunca coexistem \cite{bugden2022}.

O Algoritmo 2 demonstra uma situação em que essas regras são violadas. Inicialmente, é criada uma referência imutável para valores. Em seguida, o código tenta criar uma referência mutável para o mesmo valor. Como a referência imutável ainda está ativa, o compilador identifica o conflito entre os dois tipos de empréstimo.

\begin{lstlisting}[
    language=Rust,
    caption={Violação das regras de \emph{borrowing} em Rust},
    label={alg:borrowing-wrong}
]
fn processar_dados() {
    // Declara uma String que pode ser modificada. 
    let mut valores = String::from("Rust"); 
    
    // Cria uma referência imutável para o valor. 
    let leitura = &valores; 
    
    // Tenta criar uma referência mutável para o mesmo valor 
    // enquanto a referência imutável ainda está ativa. 
    let alteracao = &mut valores; 
    
    // Utiliza a referência imutável. 
    println!("Conteúdo: {}", leitura); 
    
    // Tenta modificar o valor por meio da referência mutável. 
    alteracao.push_str(" é seguro"); 
    
    // Exibe o valor após a alteração. 
    println!("Resultado: {}", alteracao);
}
\end{lstlisting}
\fontealgoritmo{Os autores}

O código apresentado no Algoritmo 2 é rejeitado pelo compilador antes de sua execução. O erro ocorre porque valores não pode ser emprestado de forma mutável enquanto permanece emprestado de forma imutável. Essa restrição evita que diferentes partes do programa tenham, simultaneamente, permissões incompatíveis sobre o mesmo dado.

A solução consiste em garantir que os empréstimos não se sobreponham. Para demonstrar essa situação, o Algoritmo 3 utiliza um escopo interno para limitar a duração da referência imutável. Enquanto o bloco interno está em execução, leitura pode acessar valores. Ao final desse bloco, a referência deixa de estar ativa, permitindo que uma referência mutável seja criada posteriormente.

\begin{lstlisting}[
    language=Rust,
    caption={Referências separadas por escopo, \emph{borrowing} correto},
    label={alg:borrowing-correct}
]
fn main() {
    // Declara uma String que pode ser modificada. 
    let mut valores = String::from("Rust"); 
    
    { 
        // Cria e utiliza uma referência imutável para o valor. 
        let leitura = &valores;
        println!("Conteúdo: {}", leitura);
    } // A referência imutável deixa de existir neste ponto. 
    
    // Agora é seguro criar uma referência mutável, 
    // pois não há mais uma referência imutável ativa. 
    let alteracao = &mut valores; 
    alteracao.push_str(" é seguro"); 
    
    println!("Resultado: {}", alteracao);
}
\end{lstlisting}
\fontealgoritmo{Os autores}

Diferentemente do Algoritmo 2, o Algoritmo 3 é compilado e executado normalmente, pois a referência imutável leitura deixa de existir antes da criação da referência mutável alteracao. Nesse caso, o valor pode primeiro ser consultado e, posteriormente, modificado, sem que haja sobreposição entre os empréstimos.

Essa regra é verificada pelo compilador e contribui para impedir um tipo de erro difícil de depurar chamado condição de corrida (\emph{data race}), que ocorre quando diferentes partes de um programa tentam acessar e modificar simultaneamente uma mesma região de memória sem mecanismos adequados de sincronização (Bugden; Alahmar, 2022).

Outro ponto importante é que Rust não permite manipular endereços de memória diretamente fora de contextos explicitamente marcados como \emph{unsafe}. Na prática, isso significa que todo acesso a listas e \emph{arrays} é verificado automaticamente: se o programa tentar acessar uma posição inexistente, ele encerra de forma controlada em vez de continuar operando com dados corrompidos ou expondo informações que não deveria \cite{bugden2022}. Esse tipo de proteção teria evitado o Heartbleed, uma das maiores vulnerabilidades da história da internet, que afetou uma parcela expressiva dos maiores sites HTTPS do mundo ao permitir que dados fossem lidos além dos limites permitidos da memória.

O Algoritmo 4 exemplifica esse mecanismo por meio de um vetor contendo três elementos. Como os índices válidos são 0, 1 e 2, a tentativa de acessar a posição 5 representa um acesso fora dos limites. Nesse caso, o Rust detecta a condição e interrompe a execução por meio de um \emph{panic}.

\begin{lstlisting}[
    language=Rust,
    caption={Acesso fora dos limites e encerramento por panic},
    label={alg:unsafe-array}
]
fn main() {
    // Cria um vetor contendo três elementos. 
    let numeros = vec![10, 20, 30]; 
    
    // Os índices válidos do vetor são 0, 1 e 2. 
    // O índice 5 não existe, pois está fora dos limites do vetor. 
    // O Rust detecta o acesso inválido durante a execução 
    // e encerra o programa por meio de um panic. 
    println!("{}", numeros[5]); // erro: index out of bounds
}
\end{lstlisting}
\fontealgoritmo{Os autores}

Ao executar o Algoritmo 4, a compilação é realizada normalmente, mas a execução é interrompida quando o acesso ao índice inexistente é realizado. O sistema apresenta uma mensagem indicando que o vetor possui comprimento 3, enquanto o índice solicitado é 5. Esse comportamento demonstra que o Rust realiza a verificação do limite durante a execução da operação de indexação, impedindo que o acesso inválido prossiga silenciosamente.

Por fim, Rust também elimina o conceito de ponteiro nulo, aquele valor que representa "nada"~e que em C e C++ frequentemente causa travamentos silenciosos. No lugar, Rust usa o tipo \emph{Option}, que obriga o programador a decidir explicitamente o que fazer quando um valor pode não existir \cite{bugden2022}.

O Algoritmo 5 demonstra essa abordagem por meio de uma função responsável por buscar um nome a partir de um identificador. A função retorna \texttt{Some} quando encontra um valor correspondente e \texttt{None} quando nenhum valor é encontrado. Dessa maneira, a possibilidade de ausência do resultado é representada diretamente pelo tipo de retorno da função.

\begin{lstlisting}[
    language=Rust,
    caption={Tipo \emph{Option} no lugar de ponteiro nulo},
    label={alg:option-type}
]
fn buscar_nome(id: u32) -> Option<String> {
    // Retorna um nome quando o identificador corresponde a um usuário. 
    if id == 1 { 
        Some(String::from("Olá, Mundo")) 
    } else { 
        // Indica que nenhum valor foi encontrado. 
        None 
    }
}

fn main() {
    // Realiza uma busca utilizando um identificador que não existe.
    match buscar_nome(2) {
        // Executado quando a busca retorna um valor.
        Some(nome) => println!("Encontrado: {}", nome),

        // Executado quando nenhum valor é encontrado.
        None => println!("Nenhum usuário com esse id"),
    }
}
\end{lstlisting}
\fontealgoritmo{Os autores}

No exemplo apresentado, a chamada \texttt{buscar\_nome(2)} retorna \texttt{None}, pois o identificador informado não corresponde ao valor esperado pela função. O \texttt{match} trata explicitamente essa possibilidade e exibe a mensagem correspondente. Caso a função encontrasse um usuário, retornaria \texttt{Some(nome)}, e o primeiro ramo do \texttt{match} seria executado. Dessa forma, o código não pressupõe que o valor necessariamente exista, tornando essa possibilidade explícita na estrutura do programa.

O que torna tudo isso notável é que essas verificações acontecem em tempo de compilação. Erros que em outras linguagens só aparecem quando o programa já está rodando, às vezes em produção, às vezes explorados por um hacker, em Rust são detectados antes mesmo de o programa existir \cite{bugden2022}.

Nos casos apresentados nos Algoritmos 1, 2 e 3, por exemplo, as regras de \emph{ownership} e \emph{borrowing} são verificadas durante a compilação, enquanto o acesso fora dos limites demonstrado no Algoritmo 4 é detectado durante a execução por meio de um \emph{panic}. Já o Algoritmo 5 evidencia como a possibilidade de ausência de um valor pode ser representada explicitamente pelo tipo \emph{Option}.

\section{WebAssembly: ambiente e características}

Se Rust resolve o problema de escrever código seguro e eficiente, o WebAssembly resolve um problema diferente: como fazer esse código ser executado dentro de um navegador com desempenho próximo ao nativo. É justamente essa combinação que torna as duas tecnologias complementares e justifica o foco deste trabalho.

Por muito tempo, o JavaScript foi a única linguagem executada pelos navegadores. Isso funcionou para páginas simples, mas com o crescimento de aplicações mais pesadas, editores de imagem, jogos, processamento de vídeo, suas limitações ficaram evidentes. Foi para resolver esse problema que Google, Microsoft, Mozilla e Apple se uniram para criar o WebAssembly, lançado em 2017 como padrão oficial do W3C \cite{haas2017}.

WebAssembly é um formato binário compacto de baixo nível que executa com desempenho próximo ao nativo, servindo como alvo de compilação para linguagens como Rust, C e C++. Ele não foi criado para substituir o JavaScript, mas para trabalhar ao lado dele: enquanto o JavaScript cuida da interface e das interações, o Wasm assume as partes que exigem mais processamento \cite{mdn2026}.

\subsection{Arquitetura e Funcionamento}

Um programa WebAssembly é organizado em módulos isolados, onde o código de um módulo não acessa dados de outro a menos que isso seja explicitamente permitido. O armazenamento de dados acontece em uma estrutura chamada memória linear, um vetor de bytes completamente separado do código executável e das estruturas internas do motor de execução. Isso garante que mesmo um programa com falha não consiga corromper o ambiente ao redor \cite{moric2022}.

Esse isolamento tem uma consequência prática importante: módulos de origens diferentes podem executar no mesmo processo sem interferir uns nos outros, o que viabiliza o uso do Wasm em ambientes como servidores na nuvem e dispositivos embarcados, não apenas em navegadores \cite{moric2022}.

\subsection{Segurança}

A segurança faz parte da especificação do WebAssembly desde o início, e não foi adicionada depois como um recurso extra. Um módulo Wasm não possui acesso a nenhum recurso do sistema, arquivos, redes ou dispositivos, a menos que o ambiente hospedeiro conceda esse acesso de forma explícita. Além disso, o fluxo de execução é validado antes de o programa executar por meio de um mecanismo chamado integridade do fluxo de controle (\emph{Control-Flow Integrity}), que impede que código malicioso redirecione a execução para posições não autorizadas na memória \cite{moric2022}.

Como os métodos do WebAssembly são tipados, o uso incorreto de tipos, como passar um valor do tipo errado para uma função, não é possível, o que reforça ainda mais a confiabilidade do ambiente de execução. Quando combinado com as garantias de segurança de memória do Rust, esse modelo cria uma camada dupla de proteção: o compilador Rust elimina os erros antes da geração do binário, e o WebAssembly garante que o binário gerado rode dentro de limites seguros \cite{moric2022}.

\subsection{Compilação de Rust para WebAssembly}

O processo de transformar código Rust em um módulo WebAssembly é direto e bem suportado pelo ecossistema da linguagem. O fluxo básico envolve três etapas: o código é escrito em Rust normalmente, compilado para o formato .wasm por meio de ferramentas como o wasm-pack, e então carregado e executado pelo ambiente de destino, seja um navegador, um servidor ou outro \emph{runtime} compatível.

O wasm-pack é a principal ferramenta para esse processo no ecossistema Rust. Ela cuida de compilar o código, gerar os arquivos de integração com JavaScript e empacotar tudo em um formato pronto para ser publicado ou importado por uma aplicação \emph{web}. Por baixo, o compilador Rust utiliza o \emph{backend} LLVM para gerar o binário .wasm, o mesmo \emph{backend} responsável pelas otimizações que tornam o código Rust tão eficiente em ambiente nativo \cite{moric2022}.

\begin{lstlisting}[
    language=Rust,
    caption={Função Rust exportada para WebAssembly},
    label={alg:wasm-export}
]
use wasm_bindgen::prelude::*;

#[wasm_bindgen]
pub fn fibonacci(n: u32) -> u64 {
    match n {
        0 => 0,
        1 => 1,
        _ => {
            let (mut a, mut b) = (0u64, 1u64);
            for _ in 2..=n {
                let temp = a + b;
                a = b;
                b = temp;
            }
            b
        }
    }
}
\end{lstlisting}
\fontealgoritmo{Os autores}

A macro \texttt{\#[wasm\_bindgen]} gera automaticamente o código de ligação (\emph{bindings}) necessário para que essa função seja chamada a partir do JavaScript, como mostra o Algoritmo~\ref{alg:wasm-call}.

\begin{lstlisting}[
    language=JavaScript,
    caption={Chamada da função Rust a partir do JavaScript},
    label={alg:wasm-call}
]
import init, { fibonacci } from "./pkg/meu_projeto.js";

async function executar() {
    await init(); // carrega e instancia o módulo .wasm
    const resultado = fibonacci(30);
    console.log(resultado);
}

executar();
\end{lstlisting}
\fontealgoritmo{Os autores}

Vale destacar que nem todo código Rust compila para WebAssembly sem adaptações. Funcionalidades que dependem diretamente do sistema operacional, como acesso a arquivos ou criação de \emph{threads}, não estão disponíveis no ambiente do navegador por padrão.

\section{ESTUDOS COMPARATIVOS DE DESEMPENHO}

Compreender o impacto do WebAssembly no desempenho de aplicações Rust exige fazer duas análises: primeiro, o que torna o Rust eficiente em ambiente nativo; segundo, o que acontece com esse desempenho quando o mesmo código é compilado para WebAssembly e executado dentro do navegador. As subseções a seguir percorrem essas duas camadas, formando a base comparativa que orienta a investigação empírica deste trabalho.

\subsection{Rust: Desempenho sem Runtime ou Garbage Collector}
Quando um programa é executado, ele precisa reservar espaço na memória para guardar seus dados, e também precisa liberar esse espaço quando os dados não são mais necessários. Linguagens como Java e Python fazem isso automaticamente por meio de um mecanismo chamado coletor de lixo (\emph{garbage collector}), que fica em segundo plano verificando o que pode ser descartado. Isso é conveniente, mas tem um custo: o coletor consome processamento, ocupa memória extra e pode pausar o programa em momentos imprevisíveis para fazer sua limpeza \cite{bugden2022}.

Rust não tem coletor de lixo. Em vez disso, o compilador analisa o código antes de o programa ser executado e já determina exatamente quando cada dado deve ser criado e quando deve ser descartado. Quando uma variável chega ao fim do bloco em que foi criada, sua memória é liberada automaticamente, sem nenhuma instrução extra do programador e sem nenhum custo em tempo de execução.

\begin{lstlisting}[
    language=Rust,
    caption={Memória liberada automaticamente ao fim do escopo},
    label={alg:garbage-collector}
]
fn main() {
    // 'mensagem' é criada aqui e pertence a este escopo
    let mensagem = String::from("olá, Mundo!");
    println!("{}", mensagem);
    // Ao chegar aqui, 'mensagem' sai de escopo
    // e sua memória é liberada automaticamente pelo compilador.
    // Nenhuma instrução manual é necessária.
}
\end{lstlisting}
\fontealgoritmo{Os autores}

Esse modelo elimina tanto o trabalho manual de liberar memória, que em C frequentemente gera erros, quanto a sobrecarga do coletor de lixo. O resultado é um programa que roda com eficiência comparável à de C e C++, mas sem exigir que o programador gerencie a memória na mão \cite{bugden2022}.

É esse nível de eficiência nativa que serve de referência para a comparação com o WebAssembly. Quando o mesmo código Rust é compilado para o formato .wasm e executado no navegador, parte dessas garantias de desempenho precisa ser reexaminada, pois o ambiente de execução impõe restrições e custos que não existem na compilação nativa.

\subsection{WebAssembly: Desempenho Frente ao JavaScript e à Execução Nativa}

Dentro do ambiente do navegador, o WebAssembly já representa um avanço considerável em relação ao JavaScript. Em tarefas computacionais intensas, o WebAssembly supera o JavaScript, sendo em média 1,3 vezes mais rápido que o JavaScript tradicional, seu antecessor mais próximo \cite{jangda2019}. No entanto, há uma ressalva importante: transferir dados complexos para dentro da memória do módulo tem um custo. Para passar uma \emph{string} ao Wasm, por exemplo, é necessário codificá-la, localizar espaço livre na memória linear e escrevê-la \emph{byte} a \emph{byte}, um processo que pode anular o ganho de desempenho em operações simples com texto. Quando o dado já está na memória ou o algoritmo é suficientemente complexo, o Wasm volta a superar o JavaScript por ampla margem \cite{moric2022}.

\begin{lstlisting}[
    language=Rust,
    caption={Função Rust em Wasm recebendo \emph{string} do JavaScript},
    label={alg:string-transfer}
]
use wasm_bindgen::prelude::*;

#[wasm_bindgen]
pub fn contar_vogais(texto: &str) -> u32 {
    texto
        .chars()
        .filter(|c| "aeiouAEIOU".contains(*c))
        .count() as u32
}
\end{lstlisting}
\fontealgoritmo{Os autores}

Contudo, a comparação mais relevante para este trabalho não é entre WebAssembly e JavaScript, mas entre WebAssembly e a execução nativa do próprio Rust. Nesse confronto, o cenário é diferente. Estudos mostram que aplicações compiladas para WebAssembly são executadas, em média, entre 45\% e 55\% mais devagar do que o equivalente nativo, a depender do navegador utilizado \cite{jangda2019}. Esse é exatamente o \emph{gap} que este trabalho busca caracterizar no contexto específico de aplicações Rust, contribuindo com dados concretos para orientar decisões sobre quando a portabilidade do Wasm justifica a perda de desempenho em relação à execução nativa.
